%% kapa: exam=1 round=22 year=2011 period=1 form=A subject=civil_law subject_ko=민법 seq=7 number=7
%% source: Q-Net 감정평가사 기출문제 (cst00309) artlSeq=5036790 file=extracted/1008/5036790/2011년 제22회 감정평가사 1차시험 기출문제/1교시.hwp
%% conversion: kordoc 4.24.51 -> tools/gen1.py (auto); review=proofread-sonnet; flags=none
\begin{question}{7}
대리인과 사자(使者)에 관한 설명으로 옳지 않은 것은?
\begin{choices}
  \item 사자는 의사능력자일 필요는 없다.
  \item 대리행위가 효력을 발생하기 위해서는 대리인은 행위능력자이어야 한다.
  \item 대리인은 자기가 결정한 의사를 표시하나, 사자는 본인이 결정한 의사를 표시한다.
  \item 대리의 경우, 의사표시의 하자의 유무에 관하여는 대리인을 표준으로 결정한다.
  \item 사자에 의한 의사표시는 본인이 효과의사를 결정한다.
\end{choices}
\end{question}
\answer{2}
